\section{Glossary of tools}\label{app:gloss}
\begin{description}[leftmargin=4.2cm,labelwidth=4.0cm,itemsep=3pt]
\item[Evidence (marginal likelihood)] Probability of the data under a model with its parameters integrated out; higher is better, and it already contains a penalty for flexibility (Sec.~\ref{sec:refine}).
\item[Occam term, nat] $\tfrac12p\log n$ charged per model for $p$ real coefficients and $n$ data; a nat is the unit of the logarithm of a probability ratio ($1$ nat $\approx$ a factor $2.7$).
\item[Variable projection] Fitting nonlinear parameters (poles) while the linear ones (amplitudes) are solved exactly at each step.
\item[ARD, per-branch scale] Automatic relevance determination: one prior scale per component, so an unexcited component costs no prior volume (Sec.~\ref{sec:refine}).
\item[Goodness of fit, $\log_{10}p$] A $\chi^2$ test of the residual against the known noise; $\log_{10}p<-3$ means the law is rejected.
\item[Data-equivalent] Two laws whose fits differ by less than the noise allows; the data cannot choose between them.
\item[Matrix pencil] Prony-type method: rates of an exponential sum from the shift-invariance of a Hankel matrix (Sec.~\ref{sec:modes}).
\item[Dispersion variety] $\{P(\lambda,\mu)=0\}$; points $(\lambda,i\kappa)$ are the modes of a solution (Sec.~\ref{sec:prelim}).
\item[$k$-varieties regression] $k$-means with algebraic curves as centres; alternate assignment and least-squares refit (Thm.~\ref{thm:structfit}).
\item[Cram\'er--Rao bound] Lower bound on the variance of an unbiased estimate; here for the rate of an exponential in a noisy bin (Eq.~\ref{eq:crb}).
\item[BIC / Laplace evidence] $-2\ln p(\text{data}\mid\mathcal M)\approx\mathrm{BIC}$; posterior $\propto e^{-\mathrm{BIC}/2}$ (Eq.~\ref{eq:postM}).
\item[Confluent mode] $(A+Bt)e^{\lambda t}$, produced by a repeated root (Sec.~\ref{sec:confluent}).
\item[$\gamma_N$, Friedrichs angle] Stability of splitting a sum into components; $\gamma_2=1-c_F$ (Eq.~\ref{eq:gammaN}).
\item[Dirichlet-process mixture] Mixture model with an unbounded number of components; stick-breaking weights (Sec.~\ref{sec:agg}).
\item[Pathwise conditioning] Exact GP posterior samples by correcting a prior draw with the data (Matheron's rule).
\item[Importance reweighting] Weights $p/q$ turning draws from a proposal $q$ into a posterior sample.
\item[Adjusted Rand Index] Agreement of two partitions, corrected for chance; $1$ is identical.
\item[Simulation-based calibration] Check that a posterior estimator is calibrated on simulated data (Sec.~\ref{sec:train}).
\end{description}

\section{Primer: the ideas used, at machine-learning level}\label{app:primer}
\textbf{Why $\log\hat u$ is linear in the terms.} For a constant-coefficient evolution $u_t=\sum_nA_nu$ each Fourier mode obeys $\hat u(t,\kappa)=\hat u_0(\kappa)e^{t\sum_na_n(\kappa)}$, so $\log\hat u-\log\hat u_0=t\sum_na_n(\kappa)$, linear in the unknown term coefficients. A regression on the monomials $(i\kappa)^n$ reads them off, which is the regression of Section~\ref{sec:core} on the branch $\lambda(\kappa)$.

\medskip\textbf{Why BIC and not a threshold.} A threshold on a residual decides ``is $P$ factored?'' without a scale. BIC compares models by fit and parameter count, with the noise scale estimated; two models that fit equally well are separated only by the penalty. The theory says no threshold can be right for all data (\thy{T15}).

\medskip\textbf{Why Monte Carlo and not error propagation.} A first-order propagation $\Cov(\text{out})\approx J\Cov(\text{in})J^{\top}$ requires a differentiable map; the maps here ($\Eint$, the annihilator, factor structure) jump. Sampling inputs and recording outputs gives a distribution that can have several modes.

\medskip\textbf{Relation to the original idea.} Interpolation with probabilities, an analytic map per sample, a distribution of factors and a recursion on groups are all kept. What changed: groups live in mode space and on rectangles; the number of groups comes from a DP posterior; encoders act on analytic features; every threshold is meant to be inferred.

\section{Scope: what this document contains and what it does not}\label{app:scope}
It describes the algorithm as designed in the plan and as found in the code and session reports at the time of writing. It does not reproduce the proofs (they are in the analytic paper and its ledger), does not report new experiments, and does not describe stages marked \plan{} beyond their specification.

\section{Notes for the author}\label{app:ednotes}
\begin{enumerate}[nosep]
\item This edition describes the algorithm as it stands at the handoff of 2026-11-01 (main at the T103 commit, local, not pushed). The theory-to-algorithm rows run from T1 to T103; the tags T38--T103 in Section~\ref{sec:fixes} are algorithm choices, mostly empirical, and are not listed in the appendix of theory tags.
\item The loop of Section~\ref{sec:loop} differs from the first design: it groups factors, not regions of the domain. The original stopping rule on $\bar c_F<\varepsilon_2$ is kept as specification only.
\item Numbers in the worked example come from a stand-in posterior on a grid, not from scattered data; the end-to-end figures are in Section~\ref{sec:limits}.
\item Some numerical defaults are hand-set: $\gamma_N$ floor $0.02$, SNR cut $3$, outlier cap $3s$, $P_{\mathrm{conc}}=0.9$, degree limits.
\item The test counts come from the repository's own logs and a collection of the suite ($142$ tests); the full suite was run on the working tree at the time of writing: $142$ passed ($19$ min $34$ s).
\end{enumerate}

\section{The theory tags}\label{app:tags}
The tags \thy{T$n$} refer to rows of \texttt{analysis/THEORY\_TO\_ALGO.md}, which links each choice of the algorithm to a result of the analytic work. ``Proved'' here means that the analytic work records a complete proof (status \emph{DIMOSTRATO}) or a published theorem with a checked source.
\begin{center}\small
\begin{tabular}{@{}L{1.0cm}L{12.0cm}@{}}
\toprule
Tag & Result, in one sentence\\\midrule
T1 & The solution splits as $\sum_j u_j$, $u_j\in\ker Q_j$, on rectangles and, more generally, simply connected regions; it can fail when the region has holes.\\
T2 & For evolution factors, the size $\delta(P)$ of the ambiguity of the split is the degree of the discriminant; the projector for factor $j$ has poles only where its own branch collides with another.\\
T5 & $\gamma_N$ (stability of recovering components from their sum) equals $1-c_F$ for two factors, and must be computed from the full block Gram matrix when $N\ge3$.\\
T6 & No regularisation, linear or not, recovers components better than the error factor $\gamma_N^{-1/2}$; none removes the interaction term.\\
T7 & Sparse or $\ell^1$ separation of components acts only on the finite ambiguity space; the minimum-norm split is exact and simpler.\\
T8 & No stable linear encoder makes the interaction term vanish; terms are read linearly from the \emph{log}-symbol.\\
T9 & Reading terms from spectral points is a Vandermonde inversion whose condition number blows up when the wavenumbers cluster: use a spread band.\\
T10 & \textbf{Open}: the conditioning of the Hankel/Prony step (the matrix pencil) under noise.\\
T11 & Splitting loses one derivative: the components are one order rougher than the field.\\
T12 & The interaction term and the annihilator do not depend continuously on the data in Sobolev norms.\\
T13 & Identifiability needs enough distinct modes per factor; data of exponential-polynomial type never identify the annihilator itself, only a law relative to a class.\\
T14 & Under noise $\varepsilon$ the worst-case rate for coefficient errors is of order $\varepsilon^{1/3}$.\\
T15 & There is no data-free noise threshold: the right fitting order can be found, but the threshold cannot be estimated from the data alone.\\
T17 & A fit is never certifiable from finite data: outputs are labels (fit-compatible, misfit, id-ker, undecided), never ``identified''.\\
T25 & Ehrenpreis--Palamodov: every solution on a convex region is a superposition of exponentials on the dispersion variety.\\
\bottomrule
\end{tabular}
\end{center}
